\chapter{Introduction and Problem Statement}

\section{Project mission}

The project is a statistical screening workflow for the Bourse des Valeurs
Mobilières de Tunis (BVMT), the Tunis Stock Exchange. It processes OHLCV data,
meaning open, high, low, close, and traded volume for each market observation.
It also uses stored headlines from ilboursa, a Tunisian financial-news website,
as contextual evidence.

The mission was to organize a large history of market observations into a
reviewable queue. A z-score is the observed value minus a reference mean,
divided by the reference standard deviation. In this project, the reference is
local to one symbol and uses earlier observations. Separate z-scores describe
volume and close-to-close return. The system then records nearby news, applies
a narrow watchlist rule, adds broad-market context, and presents the results in
a browser dashboard.

The central distinction is between a screening signal and an explanation. A
large statistical score states that a stored observation differs from a
symbol's earlier behavior. It does not identify the cause or establish intent.
A nearby headline can help an analyst understand context, but proximity alone
does not prove causation. A missing headline means only that the stored
same-symbol news file contains no match inside the defined interval.

\section{Why the problem is difficult}

Market records contain legitimate extremes. Thin trading can produce a large
return after a long interval. A block transaction can produce a volume spike.
An earnings announcement, dividend, capital increase, suspension, governance
decision, or broad-market movement can coincide with an abrupt change. One
threshold therefore cannot turn an unusual observation into a legal or
economic conclusion.

Sparse series also create technical problems. A close-to-close percentage
change calculated across a long interruption is not comparable with an
ordinary session return. Stable series can have a standard deviation close to
zero, which makes a ratio unstable. News can be missing, can arrive after a
market movement, or can use language outside a keyword dictionary. A useful
workflow has to record these cases explicitly rather than hide them.

The workflow separates these concerns: market data produces the signal,
quality checks record source issues, news adds context, and the watchlist
reduces the queue without changing the detector. Evaluation uses simple
baselines, while the dashboard exposes the stored evidence read-only.

\section{Objectives and scope}

The system provides the following functions:
\begin{itemize}
  \setlength{\itemsep}{0pt}
  \setlength{\parskip}{0pt}
  \item combine protected market files into an 85-series
  analytical dataset;
  \item calculate transparent prior-only volume and return scores;
  \item detect sustained reductions in trading activity through a separate
  decay rule;
  \item attach same-symbol news and deterministic headline categories;
  \item construct and refine a small review queue;
  \item evaluate the reference detector on source-backed events without
  changing its parameters;
  \item serve the seven-view React Terminal through the read-only Flask API;
  \item preserve generated outputs and evaluation hashes for verification.
\end{itemize}

The intended user is a human analyst. The interfaces support inspection of a
flag, its market history, stored news context, broad-market context, and the
rule that admitted or excluded it from the watchlist. Optional generative
triage is separate from the evidence and is not required for the dashboard.

\section{Responsible interpretation}

The project uses statistical terms precisely. A \emph{flag} is a row that
crosses the volume threshold, the return threshold, or both. A \emph{news
match} is a same-symbol record within the configured interval. A
\emph{watchlist row} satisfies the explicit screening rule or belongs to the
unexplained decay route. A \emph{hit} is a flag inside a pre-defined event
window. None of these terms is a finding of market abuse.

The event register is small: seven positive rows have dated source URLs.
The evaluation first excludes TSI because it is not a listed trading series;
its source metric therefore counts six. Only five of those six have
primary-window trading coverage. One of those five is held out. TINV,
CGF, and UBCI retain unsupported original claims or dates; TSI has a dated
source but no listed trading series. The resulting counts describe defined
examples and cannot estimate general accuracy. The report therefore values
the screening, news-context, watchlist, and Terminal workflow for human review.
